\section{Correctness and Validation}

The implementation satisfies a small set of \textbf{invariants}:
\begin{itemize}
    \item every \textbf{match} produces exactly one result message;
    \item every match has exactly one \textbf{winner};
    \item every \textbf{round} consumes an even number of players;
    \item every round halves the player count;
    \item the \textbf{championship} stops when one player remains;
    \item result ordering is stable even when \textbf{goroutines} complete out of order;
    \item invalid toss functions and invalid players are reported as \textbf{errors} instead of producing undefined behavior.
\end{itemize}

The \textbf{round barrier} is the key safety property.
Because the coordinator waits for exactly one outcome per match, the next round cannot start before the current round
is complete.
Because outcomes are stored by match index, the bracket order remains stable independently of scheduling.

Error handling is also explicit.
If the toss function is `nil`, or if it returns an invalid side, the error is propagated after all workers have been
accounted for.
This keeps the failure path well defined and prevents the coordinator from advancing with incomplete information.

The \textbf{test suite} covers player validation, coin-side parsing, single-match behavior, round execution, full championship
progression, and CLI input parsing.
Tests covering \textbf{concurrency} verify that matches in the same round start independently, that the coordinator waits for
every result, that result ordering is stable, and that invalid toss functions do not deadlock the round.
The CLI tests also check the formatted championship output for an eight-player tournament.
The championship tests additionally check the one-player case, the progression of winners across rounds, and the exact
number of rounds and matches produced by the bracket.

This gives coverage for both the functional specification and the concurrency properties required by the assignment
statement, including the absence of \textbf{shared-memory coordination} in the tournament logic.
